\subsection{The numerical realization}
\label{subsec:numerical-realization}

For $X,Y\in\mathsf O$ and $e\in\mathcal I(X,Y)$, define the partial function
$\varphi_e:|X|\rightharpoonup|Y|$ by
\[
 \varphi_e(x)=y\quad\Longleftrightarrow\quad
 \N\models G_e\bigl(\gamma_X(x),\gamma_Y(y)\bigr).
\]
Let $\mathrm{Rec}$ have the terms in $\mathsf O$ as objects and
these partial functions as morphisms, with extensional equality.
The numerical indexing represents every partial recursive function,
so this is a presentation of the recursion category of partial recursive
functions on these objects.
Its restriction idempotents are partial identities on c.e.\ subsets;
the complemented ones are partial identities on recursive subsets.

Write $\mathcal C=\mathcal C_{\mathrm{rec}}$,
$\mathscr S_{\mathrm{rec}}=\Split(\mathrm{Rec})$, and
$\mathcal B_{\mathrm{rec}}=\Tot(\mathscr S_{\mathrm{rec}})$.
An object $C$ of $\mathcal C$ is a pair $(m,C_0)$, where
$m\ge0$ and $C_0\subseteq\N^m$ is computably enumerable.
Write $|C|=C_0$ for its carrier. A morphism $f:C\to D$ is
specified by a function $|f|:|C|\to|D|$ computed by a partial
recursive program on tuple codes, defined on every input from
$|C|$. Equality is equality of these functions. In computations,
$f(x)$ abbreviates $|f|(x)$. Identities and composition are
identities and composition of the represented functions.

For each $m\geq0$, choose an injective recursive tuple coding
$\kappa_m:\N^m\to\N$ with recursive image $K_m$ and recursive
inverse $\kappa_m^{-1}:K_m\to\N^m$. Take
$\kappa_0(())=0$, $K_0=\{0\}$, and $\kappa_1=1_{\N}$;
iterated numerical pairing gives such codings for $m\geq2$.
For $C=(m,C_0)$, put $\widehat C=\kappa_m(|C|)$.


\begin{proposition}\label{prop:numerical-structure}
There is an equivalence
\[
 \operatorname{Tot}(\operatorname{Split}(\mathrm{Rec}))
       \simeq\mathcal C_{\mathrm{rec}}.
\]
The category $\mathcal C_{\mathrm{rec}}$ is regular, with finite
coproducts that are disjoint and stable under pullback. Under the
displayed equivalence, every monomorphism of $\mathcal C_{\mathrm{rec}}$
corresponds to a restriction monomorphism in
$\operatorname{Split}(\mathrm{Rec})$, and
\[
 \operatorname{Par}(\mathcal C_{\mathrm{rec}},\operatorname{Mono})
       \simeq\operatorname{Split}(\mathrm{Rec}).
\]
\end{proposition}
\begin{proof}
Define a functor $Q:\mathcal C\to\mathcal B_{\mathrm{rec}}$ by
\[
 Q(C)=(\omega,1_{\widehat C}),\qquad
 Q(f)=\kappa_n |f|\kappa_m^{-1}:
     (\omega,1_{\widehat C})\longrightarrow
     (\omega,1_{\widehat D}).
\]
For $f:C\to D$ with $C=(m,C_0)$ and $D=(n,D_0)$,
the displayed formula specifies its representing partial function.
The numerical program for $Q(f)$ first semidecides membership
in $\widehat C$, then decodes its input, computes $f$, and encodes
the output. Its domain is exactly $\widehat C$, so the displayed
split morphism is total. The inverse coding equations prove
$Q(1_C)=1_{Q(C)}$ and $Q(gf)=Q(g)Q(f)$.

For $(X,e)\in\mathcal B_{\mathrm{rec}}$, let
$D_e=\gamma_X(\operatorname{dom}e)\subseteq\N$, and define
$F:\mathcal B_{\mathrm{rec}}\to\mathcal C$ by
\[
 F(X,e)=(1,D_e),\qquad
 |F(f)|=\gamma_Y f\gamma_X^{-1}|_{D_e}:D_e\longrightarrow D_d
\]
for a total split morphism $f:(X,e)\to(Y,d)$, so that
$F(f):F(X,e)\to F(Y,d)$ belongs to $\mathcal C$.
The restriction equation $\bar f=e$ makes the displayed function defined
on all of $D_e$, and $f=dfe$ makes its values belong to $D_d$. Composition and identities
are the corresponding partial-function equations in the numerical
coordinates.
Tuple decoding gives isomorphisms
\[
 u_C:FQ(C)\xrightarrow{\cong}C
       \quad\text{in }\mathcal C,\qquad
       |u_C|=\kappa_m^{-1}|_{\widehat C}.
\]
The restriction of $\gamma_X^{-1}$ to $D_e$ gives isomorphisms
\[
 v_{(X,e)}:QF(X,e)\xrightarrow{\cong}(X,e)
       \quad\text{in }\mathcal B_{\mathrm{rec}}.
\]
For $f:C\to D$ in $\mathcal C$ and
$a:(X,e)\to(Y,d)$ in $\mathcal B_{\mathrm{rec}}$,
the coding equations give
\[
 f u_C=u_DFQ(f),\qquad
 a v_{(X,e)}=v_{(Y,d)}QF(a).
\]
Thus $u:FQ\Rightarrow1_{\mathcal C}$ and
$v:QF\Rightarrow1_{\mathcal B_{\mathrm{rec}}}$ are natural
isomorphisms, proving the first equivalence.

The singleton $(0,\{()\})$ is terminal in $\mathcal C$.
Products are Cartesian products of carriers, represented by
concatenated tuples in $\N^{m+n}$ for objects of dimensions $m,n$. If $f,g:C\to D$ are parallel morphisms, their
equalizer is the object of the same dimension as $C$ with carrier
$\{x\in|C|:f(x)=g(x)\}$, included in $C$. This carrier is c.e.:
enumerate $|C|$, run the two computations in parallel, and compare
their outputs. The product and equalizer universal morphisms are the
corresponding computable pairings and restrictions.

The object $(1,\varnothing)$ is initial. For $C=(m,C_0)$ and
$D=(n,D_0)$, choose their coproduct as the object of dimension one
with tagged carrier
\[
 |C+D|=\{2\kappa_m(c):c\in|C|\}
        \cup\{2\kappa_n(d)+1:d\in|D|\}\subseteq\N,
\]
with the two displayed encoding functions as injections.
Copairing tests the tag, decodes, and runs the corresponding
branch. These operations give the coproduct universal property.
For $f:X\to C+D$ in $\mathcal C$, let $X_C,X_D$ be the
objects of the same dimension as $X$ whose carriers are the
inverse images in $|X|$ of the even and odd summands. Both are c.e.
and their inclusions give an isomorphism $X_C+X_D\cong X$:
the inverse computes $f(x)$ to determine the tag and retains $x$.
The tag sets are disjoint, so these coproducts are disjoint and
stable under pullback.

For $f:C\to D$ in $\mathcal C$, define the image object
$\operatorname{Im}(f)$ to have the dimension of $D$ and carrier
$|\operatorname{Im}(f)|=|f|(|C|)\subseteq|D|$. This carrier is
c.e. by enumerating $|C|$ and dovetailing the computations of $f$. Its image factorization
is
\[
 C\xrightarrow{e_f}\operatorname{Im}(f)
       \xrightarrow{m_f}D,
\]
where $e_f$ is surjective and $m_f$ is the inclusion.
For any surjective morphism $q:C\to D$, define a morphism
$s_q:D\to C$ by preimage search. On input $y\in|D|$, enumerate
$|C|$ and dovetail the computations of $q$ until a value
$q(x)=y$ appears, then return $x$. Surjectivity ensures termination
on $|D|$, and $qs_q=1_D$.
If $h:C\to Z$ is constant on the fibers of $q$, put
$d=hs_q:D\to Z$. For $x\in|C|$, the elements $s_q(q(x))$ and
$x$ lie in the same fiber of $q$, so $dq=h$. Surjectivity gives
uniqueness of $d$. Hence $q$ is the coequalizer of its kernel pair.
Conversely, a regular epimorphism
cannot factor through a proper image monomorphism, so regular
epimorphisms are exactly the surjections. A pullback of a surjection
is surjective on its c.e. pullback carrier. This proves regularity.

Computable singleton maps distinguish any two elements of a carrier.
Consequently a monomorphism $m:B\to C$ in $\mathcal C$ is
injective; every injective morphism is monic. Enumerating $|B|$ and
searching for a preimage computes the partial inverse of $m$ on
$|m|(|B|)$. In numerical coordinates this gives
\[
 \mu_m:Q(C)\rightharpoonup Q(B)
       \quad\text{in }\mathscr S_{\mathrm{rec}},
\]
with
\[
 \mu_mQ(m)=1_{Q(B)},\qquad
 Q(m)\mu_m=\overline{\mu_m}.
\]
The right-hand side is the partial identity on the coded image
of $m$ in $Q(C)$.

Define a functor
\[
 \widetilde Q:\operatorname{Par}(\mathcal C,\operatorname{Mono})
                 \longrightarrow\mathscr S_{\mathrm{rec}}
\]
on objects by $\widetilde Q(C)=Q(C)$ and on a monic span by
\[
 \widetilde Q[\,C\xleftarrow{m}B\xrightarrow{g}D\,]
       =Q(g)\mu_m:Q(C)\rightharpoonup Q(D).
\]
This function is defined on the coded image of $m$ and sends
the code of $m(b)$ to the code of $g(b)$ for $b\in|B|$. An isomorphism of spans therefore leaves it
unchanged. For two composable spans, their pullback selects exactly
the pairs $(b,b')$ for which the first right leg equals the second
left leg. The same condition is the domain condition for composing
the two partial functions. Their values then agree, proving
preservation of composition. Identity spans give identities, and
the restriction span $(m,m)$ gives the partial identity on the coded image of $m$.

Define a quasi-inverse
\[
 \mathsf Q_{\mathrm{num}}:\mathscr S_{\mathrm{rec}}
       \longrightarrow\mathcal K_{\mathrm{rec}},
 \qquad
 \mathcal K_{\mathrm{rec}}
       =\operatorname{Par}(\mathcal C,\operatorname{Mono}).
\]
On objects, put $\mathsf Q_{\mathrm{num}}(X,e)=F(X,e)=(1,D_e)$.
For a partial morphism $f:(X,e)\rightharpoonup(Y,d)$, put
\[
 \widehat f=\gamma_Y f\gamma_X^{-1}:\N\rightharpoonup\N,
 \qquad U_f=\operatorname{dom}(\widehat f)\subseteq D_e.
\]
The equation $f=dfe$ gives $\widehat f(U_f)\subseteq D_d$.
Let $\iota_f:(1,U_f)\hookrightarrow(1,D_e)$ be the inclusion,
and let $g_f:(1,U_f)\to(1,D_d)$ be the morphism with underlying
function $\widehat f|_{U_f}$. Define
\[
 \mathsf Q_{\mathrm{num}}(f)=
 [\,(1,D_e)\xleftarrow{\iota_f}(1,U_f)
                      \xrightarrow{g_f}(1,D_d)\,].
\]
The carrier $U_f$ is c.e., and $g_f$ is computable and total on it,
so both legs are morphisms of $\mathcal C$.
The identity on $(X,e)$ has $U_e=D_e$ and $g_e=1_{F(X,e)}$,
so its image is the identity span.
For $g:(Y,d)\rightharpoonup(Z,c)$,
\[
 U_{gf}=\{z\in U_f:\widehat f(z)\in U_g\},\qquad
 \widehat{gf}(z)=\widehat g(\widehat f(z))\quad(z\in U_{gf}).
\]
The pullback composing the two spans is the object
\[
 P_{f,g}=(2,\{(z,w)\in U_f\times U_g:\widehat f(z)=w\}).
\]
Projection to the first coordinate gives an isomorphism
$P_{f,g}\to(1,U_{gf})$ in $\mathcal C$. Its inverse
$(1,U_{gf})\to P_{f,g}$ sends $z$ to $(z,\widehat f(z))$.
Both functions are computable and total on their source carriers.
This isomorphism commutes with both span legs, proving
preservation of composition.
The restriction of $f$ is the partial identity on its domain,
so $\mathsf Q_{\mathrm{num}}(\bar f)$ is the span
$(\iota_f,\iota_f)$. Thus $\mathsf Q_{\mathrm{num}}$ preserves
restrictions. On total morphisms it restricts to $F$, followed by
the canonical inclusion of $\mathcal C$ into $\mathcal K_{\mathrm{rec}}$.

Regard each $u_C:FQ(C)\to C$ as a total isomorphism in
$\mathcal K_{\mathrm{rec}}$. For a partial morphism $k:C\rightharpoonup D$
represented by $C\xleftarrow{m}B\xrightarrow{h}D$, its underlying
partial function has domain $|m|(|B|)$ and value
$|h|(|m|^{-1}(x))$ at $x$ in that domain. The inverse of $|m|$
on its image is computable by preimage search. Write $k(x)$ for
the value of this represented partial function. If a second span
$C\xleftarrow{m'}B'\xrightarrow{h'}D$ has the same domain and values,
then these spans are isomorphic. The functions $|m'|^{-1}|m|$ and $|m|^{-1}|m'|$
define inverse morphisms $w:B\to B'$ and $w':B'\to B$ in
$\mathcal C$. They satisfy $m'w=m$, $h'w=h$, $mw'=m'$, and
$hw'=h'$ by equality of the represented domains and values.
Both sides of
\[
 k u_C=u_D\mathsf Q_{\mathrm{num}}\widetilde Q(k)
       :FQ(C)\rightharpoonup D
\]
are defined at $z\in\widehat C$ exactly when
$\kappa_r^{-1}(z)$ belongs to the domain of $k$, where $C=(r,C_0)$.
There they have value $k(\kappa_r^{-1}(z))$. The span criterion
just proved gives equality of these morphisms.
For $a:(X,e)\rightharpoonup(Y,d)$, both sides of
\[
 a v_{(X,e)}
   =v_{(Y,d)}\widetilde Q\mathsf Q_{\mathrm{num}}(a)
   :QF(X,e)\rightharpoonup(Y,d)
\]
are defined at $z\in D_e$ exactly when $z\in U_a$, and their
value there is $a(\gamma_X^{-1}(z))$.
The components $u_C$ and $v_{(X,e)}$ therefore give natural
isomorphisms
\[
 u:\mathsf Q_{\mathrm{num}}\widetilde Q
             \Rightarrow1_{\mathcal K_{\mathrm{rec}}},\qquad
 v:\widetilde Q\mathsf Q_{\mathrm{num}}
             \Rightarrow1_{\mathscr S_{\mathrm{rec}}}.
\]
Thus $\widetilde Q$ and $\mathsf Q_{\mathrm{num}}$ are
quasi-inverse restriction functors.
\end{proof}

\begin{proposition}\label{prop:arithmetic-theory-realization}
Fix $d\ge0$, let $V=(1,\N)$, and let
$\Sigma_k=(1,\Sigma_{k,0})$ and $\Pi_k=(1,\Pi_{k,0})$
have recursive carriers for $0\le k\le d$. Then the
syntax and list algebras of Definition~\ref{def:categorical-syntax}
exist in $\mathcal C_{\mathrm{rec}}$ and are stable under parameters.
For every axiom subobject of the sentence object, the inference
endofunctor of Theorem~\ref{thm:first-order-theory} has an initial
algebra of finite derivations.
\end{proposition}
\begin{proof}
Choose the variable object $V=(1,\N)$ and, for a fixed $d\ge0$,
objects $\Sigma_k,\Pi_k\in\mathcal C_{\mathrm{rec}}$ with
recursive carriers $|\Sigma_k|,|\Pi_k|\subseteq\N$ for $0\le k\le d$.
Use the syntax codes and constructor tags of
Section~\ref{sec:st-parameterized-syntax}. The finite tree parser,
the coefficient tests, and numerical equality decide the carriers
of objects $\Tm,\Fm\in\mathcal C_{\mathrm{rec}}$. Each constructor is a primitive-recursive
function on its displayed source. The unique root decomposition
and the strictly smaller sizes of its children are the properties
of these codes established in Proposition~\ref{prop:st-syntax-initiality}.
Let $K:\mathcal C_{\mathrm{rec}}^2\to\mathcal C_{\mathrm{rec}}^2$
be the endofunctor of Appendix~\ref{app:first-order-syntax}, with
these coefficient objects. Its constructor algebra is
$\iota:K(\Tm,\Fm)\to(\Tm,\Fm)$.

Fix $P\in\mathcal C_{\mathrm{rec}}$. In
$(\mathcal C_{\mathrm{rec}}/P)^2$, put
\[
 \mathbf T_P=((P\times\Tm,\pi_P),(P\times\Fm,\pi_P)),
 \qquad \iota_P:K_P\mathbf T_P\to\mathbf T_P,
\]
where $\iota_P$ retains the parameter and applies the constructor
in the other coordinate. Take any $K_P$-algebra
\[
 \mathbf X=((X,\xi),(Y,\upsilon)),\qquad
 \alpha=(\alpha_1,\alpha_2):K_P\mathbf X\to\mathbf X.
\]
Here $\xi:X\to P$ and $\upsilon:Y\to P$ are morphisms of
$\mathcal C_{\mathrm{rec}}$, while
\[
 \alpha_1:K_{P,1}\mathbf X\to(X,\xi),\qquad
 \alpha_2:K_{P,2}\mathbf X\to(Y,\upsilon)
\]
are morphisms in $\mathcal C_{\mathrm{rec}}/P$.
For example, the function-symbol operation at arity $k$ is defined
on tuples $(p,s,x_1,\ldots,x_k)$ with $p\in|P|$,
$s\in|\Sigma_k|$, and $x_i\in|X|$ with $\xi(x_i)=p$
for $1\le i\le k$; its value belongs to $|X|$ and lies
over $p$. Its representing program halts on every such tuple.

On input $(p,t)\in|P|\times|\Tm|$ or
$(p,\varphi)\in|P|\times|\Fm|$, parse the tree and evaluate its
nodes in children-before-parent order. At each node, apply the
representing program for the corresponding component of $\alpha$
to its label, its children's values, and the same parameter $p$.
Induction on the number of nodes proves that the child values
have the required sort and projection to $p$. They therefore lie
in the source carrier of that component of $\alpha$; its program
halts and returns a value of the required sort over $p$.
The nullary step uses the variable or nullary-symbol operation.
Thus this algorithm is partial recursive and halts on every input
in each of the two displayed source carriers. It defines morphisms
\[
 h_T:P\times\Tm\to X,\qquad h_F:P\times\Fm\to Y
 \quad\text{in }\mathcal C_{\mathrm{rec}},
 \qquad \xi h_T=\pi_P,\quad\upsilon h_F=\pi_P.
\]
For $h=(h_T,h_F):\mathbf T_P\to\mathbf X$, decomposition at
the root gives
\[
 h\iota_P=\alpha K_P(h):K_P\mathbf T_P\to\mathbf X
 \quad\text{in }(\mathcal C_{\mathrm{rec}}/P)^2.
\]
If $h':\mathbf T_P\to\mathbf X$ satisfies that equation,
induction on the same root decomposition gives $h'=h$.
Consequently $(\mathbf T_P,\iota_P)$ is initial for every $P$.
Taking $P=1$ gives the syntax algebra, and the construction for
all $P$ proves stability under parameters.

Let $W\in\mathcal C_{\mathrm{rec}}$ be the object of dimension
one whose carrier is the recursive set of finite list codes over $|V|$, with constructor algebra
$\iota^L=[\nil,\cons]:1+V\times W\to W$.
Write $w_0=|\nil|(())\in|W|$ for the empty-list value.
For $P\in\mathcal C_{\mathrm{rec}}$, let
$L_P:\mathcal C_{\mathrm{rec}}/P\to\mathcal C_{\mathrm{rec}}/P$
be the slice version of $Z\mapsto1+V\times Z$ defined in
Appendix~\ref{app:categorical-hypotheses}.
Its constructor algebra has carrier $(P\times W,\pi_P)$ and
structure $\iota^L_P:L_P(P\times W,\pi_P)\to(P\times W,\pi_P)$.
For an object $(Z,z)$ of that slice, the canonical isomorphism
$\sigma_{P,Z}:(P\times V)\times_P Z\to V\times Z$
sends $((p,v),z_0)$ to $(v,z_0)$. Its inverse sends $(v,z_0)$
to $((z(z_0),v),z_0)$. For any algebra
$\beta:L_P(Z,z)\to(Z,z)$, use this isomorphism and write
$\beta_0:P\to Z$ and $\beta_1:V\times Z\to Z$ for its
two underlying operations; thus $z\beta_0=1_P$ and
$z\beta_1=z\pi_Z$.
The recursive evaluation of a finite list gives
$h_L:P\times W\to Z$ with
\[
 \begin{aligned}
 h_L(p,w_0)&=\beta_0(p),\\
 h_L(p,\cons(v,w))&=\beta_1(v,h_L(p,w)),\qquad
 zh_L=\pi_P.
 \end{aligned}
\]
Induction on list length proves halting on $|P|\times|W|$ and
uniqueness. Equivalently,
$h_L\iota^L_P=\beta L_P(h_L)$ as morphisms from
$L_P(P\times W,\pi_P)$ to $(Z,z)$ in the slice.
This proves parameterized list initiality.

The free-variable and substitution algorithms in
Propositions~\ref{prop:st-list-initiality} and
\ref{prop:st-substitution-realization} use these same codes.
Their primitive-recursive traversals therefore give morphisms
\[
 \FV_{\Tm}:\Tm\to W,\quad \FV_{\Fm}:\Fm\to W,\quad
 a_T:\Tm\times(V\times\Tm)\to\Tm,\quad
 b_F:\Fm\times(V\times\Tm)\to1+\Fm
\]
in $\mathcal C_{\mathrm{rec}}$.
For these same codes, Proposition~\ref{prop:st-substitution-realization}
proves the carrier and constructor equations in $\PA$.
They therefore hold in $\N$. Uniqueness of the parameterized folds
identifies these numerical morphisms with the categorical ones.
In particular,
$\Adm=b_F^{-1}(\Fm\hookrightarrow1+\Fm)$ is the admissibility
subobject, and $a_F:\Adm\to\Fm$ is the induced substitution
morphism. The sentence object is
$j:S\hookrightarrow\Fm$, where
$|S|=\{s\in|\Fm|:\FV_{\Fm}(s)=w_0\}$.

Fix an axiom subobject $i_A:A\hookrightarrow S$ in
$\mathcal C_{\mathrm{rec}}$. Retain $A$ itself as the object of axiom labels. The realized
rule data of Appendix~\ref{app:first-order-derivations} are
\[
 p_n:R_n\to\Fm^n,\qquad q_n:R_n\to\Fm\quad(0\le n\le2),
\]
and define
$H_A:\mathcal C_{\mathrm{rec}}/\Fm\to
\mathcal C_{\mathrm{rec}}/\Fm$ as in that appendix.
Use the finite derivation-tree codes of
Lemma~\ref{lem:syntactic-derivation-initiality} with these rule labels.
If $R_n$ has dimension $m_n$, store a label as a number
$r\in\widehat{R_n}$ and decode it by $\kappa_{m_n}^{-1}$
before applying $p_n$ or $q_n$. A code is valid when each stored
label has this membership and the ordered children's conclusions
are the premises of its decoded rule instance.
There are finitely many nodes. Dovetailing their membership and
computation tests semidecides validity, giving an object
$P_A^{\mathrm{all}}$ of dimension one whose carrier is the
c.e.\ set of valid derivation codes. Membership witnesses are not part of a
derivation code. Reading the root rule and computing its conclusion
gives $c_A^{\mathrm{all}}:P_A^{\mathrm{all}}\to\Fm$.
The node constructors give an algebra
\[
 \delta_A:H_A(P_A^{\mathrm{all}},c_A^{\mathrm{all}})
           \to(P_A^{\mathrm{all}},c_A^{\mathrm{all}})
 \quad\text{in }\mathcal C_{\mathrm{rec}}/\Fm.
\]

For any object $(Q,q)$ of $\mathcal C_{\mathrm{rec}}/\Fm$ and
algebra $\beta:H_A(Q,q)\to(Q,q)$ in this slice,
evaluate a valid derivation in children-before-parent order.
Its children have already been assigned values in $|Q|$ whose
images under $q$ are their conclusions. Together with the node's
decoded rule instance these values belong to the corresponding pullback
summand of $H_A(Q,q)$. The program representing $\beta$ therefore
halts at that input and returns a value in $|Q|$ with the node's
conclusion as its image under $q$. The same argument applies to
a nullary node with no child values. Induction on tree size proves
halting throughout the finite evaluation on every input in
$|P_A^{\mathrm{all}}|$. It gives a morphism
$h_A:P_A^{\mathrm{all}}\to Q$ in $\mathcal C_{\mathrm{rec}}$ with
\[
 qh_A=c_A^{\mathrm{all}},\qquad
 h_A\delta_A=\beta H_A(h_A).
\]
The second equation is in $\mathcal C_{\mathrm{rec}}/\Fm$, from
$H_A(P_A^{\mathrm{all}},c_A^{\mathrm{all}})$ to $(Q,q)$.
Root decomposition proves the equation, and induction on tree
size proves uniqueness of $h_A$. This establishes initiality.
\end{proof}

For arithmetic, choose exactly the signature labels and numerical
syntax codes used in Section~\ref{sec:st-parameterized-syntax}.
For a c.e.\ axiom presentation $T$ containing the chosen PA sentence
axioms, use the same axiom predicate as in
Section~\ref{subsec:arithmetic-internal-comparison}, with $\Gamma=T$.
Let
$c_T^{\mathrm{all}}:\mathrm{Prf}_T^{\mathrm{all}}\to\Fm$
be the full derivation conclusion of the proposition. Its sentence
restriction is the pullback in $\mathcal C_{\mathrm{rec}}$
\[
\begin{tikzcd}
 \mathrm{Prf}_T\arrow[r,hook]\arrow[d,"c_T"'] &
 \mathrm{Prf}_T^{\mathrm{all}}\arrow[d,"c_T^{\mathrm{all}}"]\\
 S\arrow[r,hook,"j"'] & \Fm.
\end{tikzcd}
\]
Thus $|\mathrm{Prf}_T|$ consists of the full derivation codes whose
conclusions belong to $|S|$. Theorem~\ref{thm:first-order-theory}
identifies the regular image $m_T:\mathcal T_T\hookrightarrow S$
of $c_T$ as the sentence part of the least deductively closed
formula subobject containing the axioms of $T$.

To compare with the syntactic category, write
$S^{\mathrm{syn}}\in\mathcal B_{\PA}$ and
$c_T^{\mathrm{syn}}:P_T^{\mathrm{syn}}\to S^{\mathrm{syn}}$
for the sentence object and conclusion morphism denoted by $S,c_\Gamma$
in Section~\ref{subsec:arithmetic-internal-comparison}, at $\Gamma=T$.
Let $\Fm^{\mathrm{syn}}\in\mathcal B_{\PA}$ be the formula object
with the syntax codes fixed above. Use the axiom inclusion of the
chosen sentence-code subset in both constructions. Let
$R_n^{\mathrm{num}}$ and $R_n^{\mathrm{syn}}$ be the resulting
rule objects in $\mathcal C_{\mathrm{rec}}$ and $\mathcal B_{\PA}$,
respectively, for $0\le n\le2$. Their rule morphisms are
\[
 \begin{aligned}
 p_n^{\mathrm{num}}&:R_n^{\mathrm{num}}\to\Fm^n,&
 q_n^{\mathrm{num}}&:R_n^{\mathrm{num}}\to\Fm,\\
 p_n^{\mathrm{syn}}&:R_n^{\mathrm{syn}}\to(\Fm^{\mathrm{syn}})^n,&
 q_n^{\mathrm{syn}}&:R_n^{\mathrm{syn}}\to\Fm^{\mathrm{syn}}.
 \end{aligned}
\]
Let $m_n$ be the dimension of $R_n^{\mathrm{num}}$.
The finite-product, coproduct, and pullback comparisons for these
rule constructions give bijections
\[
 \rho_n:|R_n^{\mathrm{num}}|\xrightarrow{\cong}
                  G_{\N}R_n^{\mathrm{syn}}\quad(0\le n\le2).
\]
They match the constructor parameters, tuple entries, and tags of
corresponding rule instances. These maps and their inverses are
computable on their carriers by the structural encodings and projections.
The rule membership conditions coincide because they use the same
syntax operations and axiom predicate. Let
$\chi_n:G_{\N}((\Fm^{\mathrm{syn}})^n)\xrightarrow{\cong}|\Fm|^n$
be the finite-product comparison, using the fixed equality
$G_{\N}\Fm^{\mathrm{syn}}=|\Fm|$. The definitions give
\[
 \chi_n(G_{\N}p_n^{\mathrm{syn}})\rho_n=|p_n^{\mathrm{num}}|,
 \qquad
 (G_{\N}q_n^{\mathrm{syn}})\rho_n=|q_n^{\mathrm{num}}|.
\]
These are equations of functions on $|R_n^{\mathrm{num}}|$.
Let $J_{\PA}:\mathcal P_{\PA}\to\mathcal B_{\PA}$ be the
predicate equivalence of Lemma~\ref{lem:syntactic-total-category}.
For each $n$, use the predicate presentation $C_n(z)$ and
isomorphism $u_n:J_{\PA}(C_n)\xrightarrow{\cong}R_n^{\mathrm{syn}}$
chosen in the derivation construction of
Lemma~\ref{lem:syntactic-derivation-initiality}. Put
\[
 C_n(\N)=\{r\in\N:\N\models C_n(r)\},\qquad
 \eta_n=G_{\N}(u_n):C_n(\N)\xrightarrow{\cong}
                                      G_{\N}R_n^{\mathrm{syn}}.
\]
Both $\eta_n$ and its inverse are computable on their carriers,
by the representing programs of $u_n$ and $u_n^{-1}$.
At a numerical derivation node of arity $n$ with stored label $r$,
replace that label by
$\eta_n^{-1}\rho_n(\kappa_{m_n}^{-1}(r))\in C_n(\N)$ and
recursively relabel its ordered children. The inverse replaces a
syntactic label $s\in C_n(\N)$ by
$\kappa_{m_n}\rho_n^{-1}(\eta_n(s))\in\widehat{R_n^{\mathrm{num}}}$.
Induction on tree size, using the two rule
comparison equations, proves that these procedures preserve validity
and conclusions and are inverse on the valid derivation carriers.
Restricting to sentence conclusions gives a computable bijection
\[
 \rho_T:|\mathrm{Prf}_T|\xrightarrow{\cong}G_{\N}P_T^{\mathrm{syn}},
 \qquad
 (G_{\N}c_T^{\mathrm{syn}})\rho_T=|c_T|.
\]
The sentence comparison is $G_{\N}S^{\mathrm{syn}}=|S|$.
In particular, the image of $c_T$ has underlying carrier
$\{s\in|S|:\N\models\operatorname{Prov}_T(s)\}$, for the
predicate fixed in that section.

Put $\mathscr S_{\mathrm{rec}}=\Split(\mathrm{Rec})$ and
transport these objects and total maps along $Q$.
Define
\[
 E=Q(V)=(\omega,1_\omega),\qquad
 S_{\mathrm{rec}}=Q(S),\qquad
 t_T=\ran(Q(c_T)):S_{\mathrm{rec}}\rightharpoonup S_{\mathrm{rec}}.
\]
The partial inverse of $Q(m_T)$ and
Lemma~\ref{lem:range-regular-image} identify $t_T$ with the partial
identity on the numerical image of $c_T$.
Negation and numeral substitution are the morphisms
\[
 \nu_{\mathrm{num}}:S\to S,\quad
    \nu_{\mathrm{num}}(s)=\operatorname{Neg}(s),\qquad
 \theta_{\mathrm{num}}:V\to S,\quad
    \theta_{\mathrm{num}}(n)=\Theta(n)
\]
in $\mathcal C_{\mathrm{rec}}$. Set
$\nu=Q(\nu_{\mathrm{num}}):S_{\mathrm{rec}}\to S_{\mathrm{rec}}$
and $\theta=Q(\theta_{\mathrm{num}}):E\to S_{\mathrm{rec}}$
in $\mathscr S_{\mathrm{rec}}$.
Because the syntax codes coincide, their underlying functions are
the standard denotations of the negation and numeral-substitution
morphisms of Lemma~\ref{lem:arithmetic-diagonal-representation}.
